\chapter{\texorpdfstring{\(L^p\)}{Lp} 空间中的逼近}
\label{chap:lp-approximation}

\(L^p\) 元素可以很不规则，但在 \(1\le p<\infty\) 时，它们能够由简单
函数、连续紧支撑函数乃至光滑紧支撑函数逼近。逼近定理使很多先在“好函数”
上证明的运算扩张到整个 \(L^p\) 空间。\(L^\infty\) 的情形不同：简单函数
仍可逼近单个有界函数，但空间通常不可分，连续函数也不再对本质上确界范数
稠密。

\section{简单函数逼近}

\begin{theorem}
\label{thm:simple-dense-lp}
设 \(0<p<\infty\)。对每个 \(f\in L^p(\mu)\)，存在可测简单函数
\(\varphi_n\)，使
\[
 \|\varphi_n-f\|_p\longrightarrow0.
\]
若 \(f\ge0\)，可取 \(0\le\varphi_n\uparrow f\) 几乎处处。
\end{theorem}

\begin{proof}
对 \(f\ge0\)，取第 \ref{chap:nonnegative-integral} 章的简单函数
\(0\le\varphi_n\uparrow f\)。则
\[
 0\le|f-\varphi_n|^p\le f^p,\qquad
 |f-\varphi_n|^p\to0.
\]
控制收敛定理给 \(\|\varphi_n-f\|_p\to0\)。一般实值 \(f\) 分别逼近
\(f^+,f^-\) 后相减。
\end{proof}

若测度空间 \(\sigma\)-有限，还可要求简单函数只支撑在有限测度集上。取
\(X_k\uparrow X\)、\(\mu(X_k)<\infty\)，先用
\(f\mathbf1_{X_k}\) 截断空间，再截断函数值并作简单逼近：
\[
 \|f-f\mathbf1_{X_k}\|_p^p
 =\int_{X\setminus X_k}|f|^p\longrightarrow0.
\]

\begin{corollary}
\label{cor:indicator-span-dense}
有限测度可测集的指示函数的有限线性组合在 \(L^p\) 中稠密
\((0<p<\infty)\)。
\end{corollary}

\section{连续紧支撑函数}

本节取 \(\mathbb R^n\) 上的 Lebesgue 测度。

\begin{lemma}
\label{lem:indicator-compact-open-approx}
若 \(E\subset\mathbb R^n\) 可测且 \(m(E)<\infty\)，则对每个
\(\varepsilon>0\)，存在 \(g\in C_c(\mathbb R^n)\)，满足
\[
 0\le g\le1,\qquad
 \|\mathbf1_E-g\|_p<\varepsilon
\]
对任意预先固定的 \(0<p<\infty\) 成立。
\end{lemma}

\begin{proof}
令 \(\eta=\varepsilon^p\)。由内正则性取紧集 \(K\subset E\)，使
\(m(E\setminus K)<\eta/2\)。若 \(K=\varnothing\)，这个估计已经说明
取 \(g=0\) 即可。以下设 \(K\ne\varnothing\)。由外正则性以及 \(K\) 的
有界性，可取有界开集 \(U\supset K\)，使
\[
 m(U\setminus K)<\eta/2.
\]
令
\[
 g(x)=\frac{d(x,U^c)}{d(x,K)+d(x,U^c)}.
\]
因 \(K\subset U\)，分母处处为正；\(g\) 连续，等于 \(1\) 于 \(K\)，
等于 \(0\) 于 \(U^c\)，且支撑包含于 \(\overline U\)，所以
\(g\in C_c\)。差 \(|\mathbf1_E-g|\) 不超过 \(1\)，且只可能在
\((E\setminus K)\cup(U\setminus K)\) 上非零，故
\[
 \|\mathbf1_E-g\|_p^p
 \le m(E\setminus K)+m(U\setminus K)<\eta=\varepsilon^p.
\]
\end{proof}

\begin{theorem}
\label{thm:cc-dense-lp}
对 \(1\le p<\infty\)，
\[
 \overline{C_c(\mathbb R^n)}^{\,\|\cdot\|_p}
 =L^p(\mathbb R^n).
\]
该结论对 \(0<p<1\) 也在度量
\(\int|f-g|^p\) 下成立。
\end{theorem}

\begin{proof}
先以有限测度支撑的简单函数逼近 \(f\)。每个这样的简单函数是有限个
有限测度集合指示函数的线性组合。由引理逐项用 \(C_c\) 函数逼近，并用
Minkowski 不等式控制有限和，即得 \(p\ge1\) 的结论。\(0<p<1\) 时使用
\(|\sum_ja_j|^p\le\sum_j|a_j|^p\)。
\end{proof}

\section{可分性}

\begin{theorem}
\label{thm:lp-separable}
对 \(0<p<\infty\)，\(L^p(\mathbb R^n)\) 可分。
\end{theorem}

\begin{proof}
考虑所有如下函数：支撑在有理端点有界矩形的有限并上，在有限个有理端点
半开矩形上取有理值。这样的函数只有可数多个。任意
\(C_c(\mathbb R^n)\) 函数在某个大有理立方体内一致连续，可用足够细的
有理网格和有理函数值作一致逼近；支撑所在立方体测度有限，所以一致误差
控制 \(L^p\) 误差。结合 \(C_c\) 的稠密性，该可数族在 \(L^p\) 中稠密。
\end{proof}

\begin{proposition}
\label{prop:linfty-nonseparable}
\(L^\infty(0,1)\) 不可分。
\end{proposition}

\begin{proof}
族
\[
 \{\mathbf1_{(0,t)}:0<t<1\}
\]
不可数。若 \(s\ne t\)，则两函数之差在正测区间上绝对值为 \(1\)，故
\[
 \|\mathbf1_{(0,s)}-\mathbf1_{(0,t)}\|_\infty=1.
\]
若存在可数稠密集，则半径 \(1/3\) 的球必须覆盖该族；但每个这样的球最多
包含族中一个元素，矛盾。
\end{proof}

不可分性不是说每个 \(L^\infty\) 函数都不能由简单函数逼近。事实上对
本质有界函数按值域等分，仍能在本质上确界范数中得到简单函数逼近。失败的
是存在一个对所有函数都稠密的可数简单函数族。

\section{平移与磨光}

取 \(\rho\in C_c^\infty(\mathbb R^n)\)，满足
\(\rho\ge0\)、\(\int\rho=1\)，并令
\[
 \rho_\varepsilon(x)=\varepsilon^{-n}\rho(x/\varepsilon).
\]
这样的函数族称为磨光核。对 \(f\in L^p\)，定义
\[
 (\rho_\varepsilon*f)(x)
 =\int_{\mathbb R^n}\rho_\varepsilon(y)f(x-y)\,dy.
\]

\begin{proposition}[平移连续性]
\label{prop:translation-continuity-lp}
若 \(1\le p<\infty\)，则对
\(\tau_hf(x)=f(x-h)\)，有
\[
 \|\tau_hf-f\|_p\longrightarrow0\qquad(h\to0).
\]
\end{proposition}

\begin{proof}
先设 \(f\in C_c\)。一致连续性给出点态一致小量，而所有小平移的支撑都
落在某个固定有限测度紧邻域中，因此
\(\|\tau_hf-f\|_p\to0\)。对一般 \(f\in L^p\)，取
\(g\in C_c\) 使 \(\|f-g\|_p<\varepsilon\)。平移保持 Lebesgue 测度，
故保持 \(L^p\) 范数，于是
\[
 \|\tau_hf-f\|_p
 \le2\|f-g\|_p+\|\tau_hg-g\|_p.
\]
先令 \(h\to0\)，再令 \(\varepsilon\downarrow0\)。
\end{proof}

\begin{theorem}[磨光逼近]
\label{thm:mollifier-approximation}
若 \(1\le p<\infty\)，\(f\in L^p(\mathbb R^n)\)，则
\[
 \rho_\varepsilon*f\in C^\infty(\mathbb R^n),\qquad
 \|\rho_\varepsilon*f-f\|_p\longrightarrow0.
\]
若先把 \(f\) 截断为紧支撑，再磨光，则可取逼近函数属于
\(C_c^\infty(\mathbb R^n)\)。因此
\[
 \overline{C_c^\infty(\mathbb R^n)}^{\,\|\cdot\|_p}
 =L^p(\mathbb R^n).
\]
\end{theorem}

\begin{proof}
卷积的光滑性由把导数移到光滑紧支撑核上得到：
\[
 \partial^\alpha(\rho_\varepsilon*f)
 =(\partial^\alpha\rho_\varepsilon)*f.
\]
局部积分可由 Hölder 不等式控制。由于 \(\rho_\varepsilon\) 是概率密度，
Jensen 不等式给出
\[
 \left|\int\rho_\varepsilon(y)
   \bigl(f(x-y)-f(x)\bigr)\,dy\right|^p
 \le\int\rho_\varepsilon(y)|f(x-y)-f(x)|^p\,dy.
\]
对 \(x\) 积分并用 Tonelli 定理，得到
\[
 \begin{split}
 \|\rho_\varepsilon*f-f\|_p^p
 &\le\int\rho_\varepsilon(y)\|\tau_yf-f\|_p^p\,dy\\
 &\le\sup_{|y|\le C\varepsilon}\|\tau_yf-f\|_p^p
 \longrightarrow0,
 \end{split}
\]
其中 \(C\) 使 \(\operatorname{supp}\rho\subset B(0,C)\)。若
\(f\) 有紧支撑，则卷积支撑包含于两个紧支撑的 Minkowski 和中，仍为紧
支撑。先作空间截断，再磨光，即得最后结论。
\end{proof}

\section{从稠密子空间延拓算子}

\begin{proposition}[有界算子延拓]
\label{prop:bounded-operator-extension-dense}
设 \(D\) 是赋范空间 \(X\) 的稠密线性子空间，\(Y\) 是 Banach 空间。
若线性算子 \(T:D\to Y\) 满足
\[
 \|Tf\|_Y\le C\|f\|_X\qquad(f\in D),
\]
则存在唯一有界线性算子
\(\widetilde T:X\to Y\) 延拓 \(T\)，且
\(\|\widetilde T\|\le C\)。
\end{proposition}

\begin{proof}
对 \(f\in X\)，取 \(f_n\in D\) 使 \(f_n\to f\)。估计
\[
 \|Tf_n-Tf_m\|_Y\le C\|f_n-f_m\|_X
\]
说明 \((Tf_n)\) 在完备空间 \(Y\) 中收敛。定义
\(\widetilde Tf=\lim_nTf_n\)。同一估计说明定义与逼近列无关，并给出
线性、有界性。若两个连续延拓在稠密集 \(D\) 上相同，则在 \(X\) 上相同。
\end{proof}

这一命题解释了为什么可以先在 \(C_c^\infty\) 上定义积分算子、Fourier
变换或弱导数估计：只要范数估计与所选光滑逼近无关，算子就有唯一连续延拓。

\section*{习题}
\addcontentsline{toc}{section}{习题}
\subsection*{A组　稠密性}
\begin{enumerate}[label=\textbf{\thechapter.\arabic*.}]
\exerciseitem{simple-dense-complex}
把简单函数稠密性推广到复值 \(L^p\) 函数。

\exerciseitem{finite-support-simple-dense}
在 \(\sigma\)-有限测度空间上，完整证明有限测度支撑简单函数在
\(L^p\) 中稠密。

\exerciseitem{compact-open-distance-function}
验证引理 \ref{lem:indicator-compact-open-approx} 中距离函数公式的连续性、
取值和紧支撑性质。

\exerciseitem{cc-density-error-budget}
设 \(\varphi=\sum_{j=1}^ma_j\mathbf1_{E_j}\)。给定
\(\varepsilon>0\)，明确分配每项误差以构造
\(g\in C_c\) 使 \(\|\varphi-g\|_p<\varepsilon\)。

\end{enumerate}
\subsection*{B组　可分性与平移}
\begin{enumerate}[label=\textbf{\thechapter.\arabic*.},resume]
\exerciseitem{lp-separable-rational-grid}
把定理 \ref{thm:lp-separable} 中的可数稠密族写成精确定义，并证明它
确实可数。

\exerciseitem{linfty-nonseparable-balls}
补全 \(L^\infty(0,1)\) 不可分证明中的球覆盖论证。

\exerciseitem{translation-isometry}
证明 Lebesgue 测度下
\(\|\tau_hf\|_p=\|f\|_p\)，并说明一般测度上为何未必成立。

\exerciseitem{translation-discontinuous-linfty}
给出 \(f\in L^\infty(\mathbb R)\)，使
\(\|\tau_hf-f\|_\infty\not\to0\) 当 \(h\to0\)。

\end{enumerate}
\subsection*{C组　磨光与算子延拓}
\begin{enumerate}[label=\textbf{\thechapter.\arabic*.},resume]
\exerciseitem{mollifier-norm-bound}
证明
\[
 \|\rho_\varepsilon*f\|_p\le\|f\|_p
\]
对 \(1\le p\le\infty\) 成立。

\exerciseitem{mollifier-support}
若 \(\operatorname{supp}f\subset K\)、
\(\operatorname{supp}\rho_\varepsilon\subset B(0,\varepsilon)\)，证明
\[
 \operatorname{supp}(\rho_\varepsilon*f)
 \subset K+\overline{B(0,\varepsilon)}.
\]

\optionalexerciseitem{density-open-set}
设 \(\Omega\subset\mathbb R^n\) 开。证明
\(C_c^\infty(\Omega)\) 在 \(L^p(\Omega)\) 中稠密
\((1\le p<\infty)\)。

\projectexerciseitem{operator-extension-application}
先在有限支撑序列上定义离散卷积
\((a*b)_n=\sum_ka_{n-k}b_k\)，证明合适的 Young 范数估计，再用稠密延拓
把算子定义到相应的 \(\ell^p\) 空间。
\end{enumerate}
